\section{预训练数据与伦理考量}

\subsection{预训练数据的来源}

\begin{figure}[H]
\centering
\includegraphics[width=0.95\textwidth]{slides-images/slide-025.jpg}
\caption{预训练数据的来源：左图为 The Pile 数据集的组成（学术论文、互联网爬取、书籍、代码等）；右表为各模型使用的训练数据，GPT-3.5+ 的数据已不再公开\protect\footnotemark}
\end{figure}
\footnotetext{Slides 第26页。}

主要的预训练数据来源包括：

\begin{itemize}
    \item \textbf{CommonCrawl}：互联网爬取数据（量最大，质量参差不齐）
    \item \textbf{Wikipedia}：高质量百科全书文本
    \item \textbf{BookCorpus / Books}：书籍语料
    \item \textbf{WebText / OpenWebText}：从 Reddit 高赞链接中收集的网页
    \item \textbf{学术论文}（ArXiv、PubMed）
    \item \textbf{代码}（GitHub、StackExchange）
\end{itemize}

\subsection{预训练学到了什么——以及它的阴暗面}

\begin{figure}[H]
\centering
\includegraphics[width=0.95\textwidth]{slides-images/slide-060.jpg}
\caption{预训练教会了模型什么：从事实知识、句法、共指、语义、情感到推理——但模型也会学到并放大训练数据中的种族歧视、性别偏见等有害内容\protect\footnotemark}
\end{figure}
\footnotetext{Slides 第61页。}

\begin{warningbox}{预训练的伦理风险}
Chris Manning 特别强调：模型在学习有用知识的同时，也会``learn and exacerbate racism and sexism, all manner of biases''。这些偏见来自训练数据中的不平衡和有害内容，而模型会忠实地复制甚至放大这些偏见。

这引出了关于预训练的重要伦理问题：
\begin{itemize}
    \item 训练数据的筛选和清洗至关重要
    \item 模型的输出不应被盲目信任
    \item 部署前需要仔细评估模型在不同群体上的表现差异
    \item 大模型的能力边界和失败模式仍不完全清楚
\end{itemize}
\end{warningbox}

\subsection{本章小结}

预训练数据的规模和多样性是模型能力的基础，但数据质量和伦理问题同样不可忽视。随着模型越来越大、能力越来越强，理解它们的工作原理和失败模式变得愈发重要——这不仅是为了用好它们，更是为了负责任地部署它们。

%% ========== 总结与延伸 ==========

